%% SPDX-License-Identifier: CC-BY-4.0
\section{The reader from inside: two levels of composition}\label{sec:two-levels}

Equation \eqref{eq:I} is an extensional presentation of a finite reading. It need not be understood as a view from nowhere in which the whole graph is simultaneously present to an external collector. An internal reader may traverse the declared finite edge list one edge at a time, obtain the corresponding contribution, and update a finite accumulator. The list representation supplies an implementation order; the theorem concerns the value delivered by the declared finite computation, not an epistemic privilege attached to that order.

This raises a genuine typing question. Three states are easily collapsed in prose:
\[
\iota\quad\text{(the accumulator has received no input)},
\qquad
\bot\quad\text{(a reading failed to resolve)},
\qquad
\mathrm{Resolved}(\zeroQ)\quad\text{(a reading succeeded and retained no difference)}.
\]
They are not interchangeable. The first is a process state, the second a failure verdict, and the third a determinate quantitative verdict.

The algebraic source of the confusion is now machine-checked in \verb|ReaderTwoLevels.v|. For one operation, if a single element is both a two-sided identity and a two-sided absorber, every two carrier elements are equal. But the restriction is explicitly \emph{for one operation}. A finite reader performs at least two different compositions:

\begin{center}
\small
\begin{tabular}{@{}p{0.22\linewidth}p{0.40\linewidth}p{0.27\linewidth}@{}}
\toprule
operation & level & boundary behaviour \\
\midrule
accumulation & combine contributions within one reading & an initial state is neutral \\
pipeline sequencing & compose completed reading stages & an unresolved stage is absorbing \\
\bottomrule
\end{tabular}
\end{center}

A common carrier element can therefore be neutral for the first operation and absorbing for the second without collapse; the formal file provides an explicit witness. That witness establishes algebraic compatibility only. It does not by itself show that ``no accumulated contribution yet'' and ``the reader failed'' have the same meaning. The typed refinement keeps them separate: \verb|Init| is the unit of accumulation, \verb|AccFailed| is absorbing within a failed accumulation, and \verb|Unresolved| is absorbing under pipeline sequencing. A resolved zero remains a different constructor and behaves differently: it allows a later resolved stage to proceed, whereas an unresolved verdict stops the pipeline.

The empty-input case then becomes visible as a reading-contract decision rather than a theorem of accumulation alone. A \emph{total} contract may finalise an untouched accumulator as \(\mathrm{Resolved}(\zeroQ)\): the traversal completed successfully and accumulated no contribution. A \emph{strict} contract may finalise the same untouched process state as \(\bot\): no observation was made, so no quantitative verdict is licensed. Both finalisers are formally defined, and the development proves that they disagree on the empty case.

This qualification blocks an overclaim. Accumulation as such does not force zero. Zero is forced only after additional choices: the process is totalised on an empty trace, the result must inhabit the same quantitative codomain as ordinary readings, and the chosen accumulation law has arithmetic zero as its identity. None of this weakens the kernel theorem. Once the operator legitimately returns \(\zeroQ\), Theorems \ref{thm:edge} and \ref{thm:component} still determine exactly what that returned value certifies. The new separation says only that a process state must not be smuggled into the verdict type and then mistaken for such a return.
